\documentclass{article}
\usepackage{amsmath, amssymb, amsthm}
\usepackage{hyperref}
\usepackage{geometry}
\geometry{margin=1in}

\title{Erd\H{o}s Problem \#1101}
\date{Last updated: 2025-10-19}
\author{Source: erdosproblems.com}

\begin{document}
\maketitle

\noindent\textbf{Status:} OPEN \quad \textbf{No prize}

\noindent\textbf{Tags:} number theory

\medskip

\noindent\textbf{Problem Statement:}

If $u=\{u_1&lt;u_2&lt;\cdots\}$ is a sequence of integers such that $(u_i,u_j)=1$ for all $i\neq j$ and $\sum \frac{1}{u_i}&lt;\infty$ then let $\{a_1&lt;a_2&lt;\cdots\}$ be the sequence of integers which are not divisible by any of the $u_i$. For any $x$ define $t_x$ by\[u_1\cdots u_{t_x}\leq x&lt; u_1\cdots u_{t_x}u_{t_x+1}.\]We call such a sequence $u_i$ good if, for all $\epsilon&gt;0$, if $x$ is sufficiently large then\[\max_{a_k&#60;x} (a_{k+1}-a_k) &#60; (1+\epsilon)t_x \prod_{i}\left(1-\frac{1}{u_i}\right)^{-1}.\]Is there a good sequence such that $u_n&#60; n^{O(1)}$? Is there a good sequence such that $u_n\leq e^{o(n)}$?




Erd\H{o}s \cite{Er81h} believed the answer to the first question is no and the second question is yes. He proved the existence of some good sequence (in which all the $u_i$ are primes). An easy sieve argument proves that we always have, for any sequence $u$ with those properties,\[\max_{a_k&lt;x} (a_{k+1}-a_k)&gt; (1+o(1))t_x \prod_{i}\left(1-\frac{1}{u_i}\right)^{-1}.\]The strong form of [208] is asking whether if $u_i=p_i^2$, the sequence of prime squares, is good.




References

[Er81h] Erd\H{o}s, P., Some problems and results on additive and multiplicative number theory . Analytic number theory (Philadelphia, Pa., 1980) (1981), 171-182.


\bigskip
\noindent\small{Source: \url{https://www.erdosproblems.com/1101}}
\end{document}
